\newpage
\phantomsection
\label{prf:multiplicative-cancellation}

\begin{remark*}[Return]
\hyperref[thm:multiplicative-cancellation]{Return to Theorem}
\end{remark*}

\begin{theorem*}[Multiplicative Cancellation]
Let $F$ be a field and let $a,b,c\in F$ with $c\neq 0$.  If
\[
a\cdot c = b\cdot c,
\]
then $a = b$.
\end{theorem*}

\begin{proof}
\textbf{Professional Standard Proof.}~\\
Assume $a\cdot c=b\cdot c$ with $c\neq 0$.  Since $F$ is a field, $c$ has a
multiplicative inverse $c^{-1}$.  Multiply the given equality on the right by
$c^{-1}$:
\[
(a\cdot c)\cdot c^{-1}=(b\cdot c)\cdot c^{-1}.
\]
By associativity,
\[
a\cdot(c\cdot c^{-1})=b\cdot(c\cdot c^{-1}).
\]
Since $c\cdot c^{-1}=1$, this becomes
\[
a\cdot 1=b\cdot 1.
\]
By the multiplicative identity axiom, $a=b$.
\end{proof}

\begin{remark*}[Proof structure]
The proof multiplies both sides of the equation by the inverse of the common
nonzero factor.  Associativity moves the factor $c^{-1}$ next to $c$, the
inverse axiom collapses $c\cdot c^{-1}$ to $1$, and the identity axiom then
removes the surviving factor of $1$.  This is the multiplicative analogue of
the usual additive cancellation argument.
\end{remark*}

\begin{remark*}[Dependencies]~\\
\begin{itemize}
  \item \hyperref[def:field]{Definition~\ref*{def:field}}
        \textnormal{(multiplicative inverse; multiplicative identity; associativity of multiplication)}
\end{itemize}
\end{remark*}
